\subsection{An absolute essential-image criterion}

The criterion below is an existence statement about explicit data on
the closed locus, its infinitesimal neighbourhoods, the completion,
and the punctured completion. In particular it does not quantify over
an abelian scheme on the original base. Its two effectivity conditions
are retained because unpolarized formal algebraization and descent of
schemes cannot be taken for granted.

\begin{theorem}[C-RD13: the exact image by formal models and compatible charts]
\label{thm:C-RD13}
\status{Proved; independent root audit passed.}
Let $S$ be a locally noetherian scheme, let $T\subset V(p)$ be a closed subset,
put $U=S\setminus T$, and fix a triple $(A_U,M,\varphi)$ as in C.1.
Choose a coherent ideal $\mathcal I\subset\mathcal O_S$ containing
$p$, with $V(\mathcal I)=T$, and put
\[
 T_r=V(\mathcal I^r)\quad(r\geq1).
\]
The triple is in the essential image of $F$ if and only if there
exist the following data.

\begin{enumerate}
\item An abelian scheme $A_1/T_1$ and a specified isomorphism
$\eta_1:A_1[p^\infty]\simeq M_{T_1}$.
For $r\geq1$, let $(A_r,\eta_r)$ be the compatible Serre--Tate
liftings of this datum to $T_r$, using $M_{T_r}$. These are determined
up to their unique reduction-preserving compatible isomorphisms.

\item An affine open cover $S=\bigcup_i V_i$, where
$V_i=\Spec R_i$, and, for each $i$, an actual abelian scheme
$B_i/\widehat V_i$, with
\[
 \widehat V_i=\Spec\widehat{R_i}_{I_i},\qquad
 I_i=\mathcal I(V_i),\qquad
 \widehat U_i=\widehat V_i\times_SU,
\]
and compatible identifications
$B_i\times_{\widehat V_i}(T_r\cap V_i)\simeq A_r|_{T_r\cap V_i}$.
Thus $B_i$ is an algebraization of a specified formal object, rather
than an unspecified good model. Its formal BT identifications
uniquely give an actual isomorphism
\[
 \widehat\eta_i:B_i[p^\infty]\simeq M_{\widehat V_i}.
                                                        \tag{C-RD13.1}
\]

\item A group-scheme isomorphism
\[
 \beta_i:B_i|_{\widehat U_i}\ \simeq\ A_U|_{\widehat U_i}
                                                        \tag{C-RD13.2}
\]
whose map on BT groups is
$\varphi|_{\widehat U_i}\circ\widehat\eta_i|_{\widehat U_i}$.
The condition requires an actual abelian-scheme isomorphism on this
overlap, and not merely an isomorphism of its BT groups.

\item For each $i$, finite open covers
\[
 B_i=\bigcup_a W_{ia},\qquad
 A_U|_{U\cap V_i}=\bigcup_a Z_{ia},
                                                        \tag{C-RD13.3}
\]
where $W_{ia}\to\widehat V_i$ and $Z_{ia}\to U\cap V_i$ are
quasi-affine, with a common finite index set, and $\beta_i$ identifies
$W_{ia}|_{\widehat U_i}$ with $Z_{ia}|_{\widehat U_i}$ for every $a$.
Empty members are allowed.
\end{enumerate}
The ideal and affine cover may be chosen arbitrarily: existence for
one such choice is equivalent to existence for every choice. No
excellence assumption or finite-normalization assumption occurs.
The statement therefore applies in particular under a) and b) of C.1,
and on reduced and nonreduced bases alike.
\end{theorem}
\begin{proof}
We first justify the objects and comparisons in the statement.
Since $p\in\mathcal I$, the element $p$ is nilpotent on every $T_r$.
Each immersion $T_r\hookrightarrow T_{r+1}$ is nilpotent. The
Serre--Tate equivalence
\cite[Theorem~1.2.1]{Katz1981SerreTate} consequently constructs the
compatible liftings $A_r$, including their BT identifications, from
$(A_1,\eta_1)$. This construction works globally on the closed
schemes by the same full faithfulness and gluing. Given $B_i$ in
condition 2, for every fixed $n$ the coordinate algebras of
$B_i[p^n]$ and $M[p^n]_{\widehat V_i}$ are finite projective modules
over the complete noetherian ring $\widehat R_i$. The compatible
isomorphisms modulo $I_i^r$ therefore give a unique Hopf-algebra
isomorphism over $\widehat R_i$. Passing through all $n$ proves
(C-RD13.1), including compatibility between levels.

Necessity is direct but also checks the atlas requirement. If $A/S$
realizes the triple, take $A_1=A|_{T_1}$ and
$B_i=A\times_S\widehat V_i$. Serre--Tate identifies their formal
reductions with the stated $A_r$, and the given identification with
$A_U$ supplies $\beta_i$. Choose a finite affine open cover of
$A|_{V_i}$, which is possible because it is a noetherian
quasi-compact scheme. An affine scheme mapping to the affine $V_i$
is affine over $V_i$. Its two base changes are therefore affine,
hence quasi-affine, over $\widehat V_i$ and $U\cap V_i$. These are
the paired opens in (C-RD13.3). This proves necessity for every
choice of $\mathcal I$ and of the affine cover.

For sufficiency first work on one affine $V=\Spec R$, omitting $i$.
The completion morphism $\widehat V\to V$ is flat, its inverse image
of $T_1\cap V$ is that same scheme, and
\[
             (U\cap V)\amalg\widehat V\longrightarrow V
                                                        \tag{C-RD13.4}
\]
is faithfully flat and quasi-compact. Apply the quasi-affine gluing
equivalence of \cite[Theorem~1.1, p.565]{MoretBailly1996Descent}
to each paired open $(W_a,Z_a)$ with its restricted $\beta$.
It gives a quasi-affine $V$-scheme $X_a$ with those two prescribed
restrictions. Intersections $W_a\cap W_b$ and $Z_a\cap Z_b$ are
again relatively quasi-affine and are paired by $\beta$. The same
theorem gives $X_{ab}$ and morphisms
\[
                         X_a\longleftarrow X_{ab}\longrightarrow X_b.
\]
These morphisms are open immersions: that property can be checked
after the faithfully flat cover (C-RD13.4), where they are the
original open inclusions. By full faithfulness, their transition
isomorphisms satisfy the cocycle condition, since the condition
holds on each of the two restrictions. Gluing the $X_a$ along these
open subschemes gives a $V$-scheme $A_V$ with specified isomorphisms
\[
 A_V|_{U\cap V}\simeq A_U|_{U\cap V},\qquad
 A_V\times_V\widehat V\simeq B.
                                                        \tag{C-RD13.5}
\]
Their comparison is exactly $\beta$. Finite presentation, smoothness,
properness and geometrically connected fibres are faithfully flat
local properties on the base; they hold on (C-RD13.4), and hence
hold for $A_V/V$.

The group structure also descends without an extra effectiveness
assumption. The two additions give compatible morphisms on the two
restrictions of $A_V\times_VA_V$. The full faithfulness for schemes
in \cite[Theorem~1.2, p.565]{MoretBailly1996Descent} gives a unique
addition on $A_V$. Apply the same assertion to the identity and
inverse. The group identities hold after (C-RD13.4), hence hold on
$A_V$. Thus $A_V$ is an abelian scheme. At every finite level, the
two isomorphisms $A_V[p^n]\simeq M[p^n]$ on this cover agree on
$\widehat U$ by condition 3. The quasi-coherent gluing equivalence
\cite[Theorem~0.3, p.561]{MoretBailly1996Descent}, applied to the
finite coordinate algebras, glues them to an isomorphism over $V$.
All Hopf identities and all transitions in $n$ glue as equalities.
We have therefore realized the complete triple over $V$.

We spell out why these affine constructions glue, even though their
completed models were chosen separately. Take an affine open
$V=\Spec R$ inside $V_i\cap V_j$ and complete it along $\mathcal I$.
The base changes of $B_i,B_j$ to $\widehat V$ are proper schemes
whose formal completions are both the prescribed tower $A_r$ on
$T_r\cap V$. Full faithfulness of formal completion for proper
schemes over a complete noetherian ring gives the unique
isomorphism between these two models inducing those formal
identifications; see \cite[Theorem~5.4.1, pp.156--157]{EGAIII1}.
It preserves the group structures and their BT comparisons, since
both assertions are equalities that hold on every formal reduction.
On $\widehat V\times_SU$, it agrees with the comparison through
$A_U$: the two isomorphisms induce the same BT isomorphism, and the
map from homomorphisms of abelian schemes to homomorphisms of their
BT groups is faithful. The latter faithfulness follows on geometric
fibres from the schematic density of all $p$-power torsion, and
then on the noetherian base from unramifiedness of the abelian Hom
functor. Moret--Bailly's Theorem~1.2 now gives the unique
isomorphism between $A_{V_i}|_V$ and $A_{V_j}|_V$ with the given
open and completed restrictions. These isomorphisms agree on
smaller affine opens and satisfy the cocycle condition by the same
full faithfulness. They therefore glue the $A_{V_i}$ to an abelian
scheme on $S$, and their BT identifications glue as well.

This proves sufficiency. Combined with necessity, it proves the
independence of the chosen ideal and cover asserted in the theorem.
\end{proof}

\begin{corollary}[C-RD14: the exact elliptic criterion]
\label{cor:C-RD14}
\status{Proved; independent root audit passed.}
Suppose $A_U$ has relative dimension one and $M$ has height two. In Theorem~\ref{thm:C-RD13}, conditions 2 and 4 may be
removed. More explicitly, choose an elliptic scheme $A_1/T_1$ with
$A_1[p^\infty]\simeq M_{T_1}$. On each affine completion the
Serre--Tate tower has a unique algebraization $B_i$, and the triple
extends if and only if these $B_i$ admit the actual overlap
isomorphisms (C-RD13.2) respecting the specified BT comparisons.
There is no additional ample-class or scheme-atlas condition in
relative dimension one.
\end{corollary}
\begin{proof}
The zero sections of the elliptic schemes $A_r$ define compatible
relatively ample invertible sheaves
$\mathcal L_r=\mathcal O_{A_r}(3e_r)$. On the affine closed base
$T_1\cap V_i$ their restrictions are ample as invertible sheaves.
Grothendieck's algebraization theorem
\cite[Theorem~5.4.5, pp.157--158]{EGAIII1} gives a projective scheme $B_i$ over
$\widehat V_i$ with this formal completion and this ample line
bundle. The group operations algebraize by proper formal full
faithfulness. At a point over $V(I_i)$, all the quotients
$B_i\bmod I_i^r$ are flat over $\widehat R_i/I_i^r$. The local
formal criterion for flatness, applied to these noetherian local
rings, makes $B_i$ flat at that point. Its fibre is the smooth
fibre of the elliptic scheme $A_1$, so the finite-presentation
smoothness criterion makes $B_i$ smooth there. The complement
of the smooth locus has closed image by properness, and every
nonempty closed subset of the complete affine
base meets $V(I_i)$ because $I_i$ is in the Jacobson radical.
Geometric connectedness then follows from the finite etale
component scheme of a proper smooth group. Thus $B_i$ is elliptic.
The finite-level completeness argument supplies its prescribed BT
comparison, and formal full faithfulness gives its uniqueness.

Any group isomorphism $\beta_i$ preserves the zero sections and
therefore identifies the canonical line bundles
$\mathcal O_{B_i}(3e)$ and $\mathcal O_{A_U}(3e)$ on the overlap.
Apply \cite[Theorem~1.1, the QPROJ case, p.565]{MoretBailly1996Descent}
to these paired schemes with their relatively ample line bundles.
This directly produces a scheme on $V_i$, so no choice of a
compatible quasi-affine cover is necessary. The properness, group
structure, BT identification and global gluing follow as in the
proof of Theorem~\ref{thm:C-RD13}. Necessity remains restriction.
\end{proof}

\begin{remark}[Relation to the nilpotent correction]
The general criterion retains two distinct effectivity checks:
algebraization of the formal Serre--Tate object and a compatible
scheme atlas on the open-completion pair. They are explicit
conditions on the given triple and auxiliary data away from $S$;
neither says that the triple already extends over $S$.
Theorem~\ref{thm:C-RD12} is stronger when a reduced realization is
already specified: it proves that both checks can then be removed,
leaving exactly the additional completed-overlap isomorphism.
Corollary~\ref{cor:C-RD14} gives a second situation where both checks
disappear. The reduced counterexample C-RD11 and the square-zero
counterexample C-RD5 show that the remaining comparison cannot be
replaced by reducedness of the local rings.
\end{remark}
